\chapter*{Course Map: From R Commands to Data Insight}
\addcontentsline{toc}{chapter}{Course Map: From R Commands to Data Insight}
\markboth{Course Map}{}

The course progresses from R fundamentals to data structures, data management, visualization, and introductory analysis.

\begin{dsfigure}
\begin{tikzpicture}[
  stage/.style={draw=DSBlue,fill=DSPaleBlue,rounded corners=2mm,text width=12.4cm,minimum height=1.36cm,align=left,inner xsep=4mm,inner ysep=2.5mm,font=\scriptsize\color{DSNavy}},
  arr/.style={-{Stealth[length=2.2mm]},thick,draw=DSTeal}
]
\node[stage] (s1) at (0,5.4) {{\normalsize\bfseries\color{DSBlue}Stage 1: R orientation and language foundations}\hfill{\footnotesize Labs 1--2}\par
\smallskip Set up R/RStudio; use scripts, variables, arithmetic, logic, comparisons, and strings.};
\node[stage] (s2) at (0,3.45) {{\normalsize\bfseries\color{DSBlue}Stage 2: Core R data structures}\hfill{\footnotesize Labs 3--4}\par
\smallskip Construct, index, and manipulate vectors, lists, factors, matrices, and arrays.};
\node[stage] (s3) at (0,1.50) {{\normalsize\bfseries\color{DSBlue}Stage 3: Programming flow and data frames}\hfill{\footnotesize Labs 5--7}\par
\smallskip Apply control flow; create, import, clean, filter, sort, and export data frames.};
\node[stage] (s4) at (0,-0.45) {{\normalsize\bfseries\color{DSBlue}Stage 4: Visualization and introductory analysis}\hfill{\footnotesize Labs 8--9}\par
\smallskip Create charts, compute Pearson correlation, build heatmaps, and compare transformations and scaling methods.};
\draw[arr] (s1.south) -- (s2.north);
\draw[arr] (s2.south) -- (s3.north);
\draw[arr] (s3.south) -- (s4.north);
\end{tikzpicture}
\end{dsfigure}

\begin{keyidea}
The sequence is cumulative: Labs 7--9 assume proficiency with indexing, conditions, data frames, packages, and files.
\end{keyidea}

\begin{dstable}
\begin{tabularx}{\linewidth}{@{}p{1.4cm}p{4.1cm}X@{}}
\toprule
\rowcolor{DSPaleBlue!92}
Lab & Main topic & Practical emphasis \\
\midrule
1 & RStudio, R, variables, arithmetic & Set up R; write code; use assignment, types, arithmetic, and input. \\
2 & Logic, relations, strings & Evaluate Boolean conditions and manipulate text. \\
3 & Vectors and lists & Store, index, and manipulate homogeneous and mixed collections. \\
4 & Factors, matrices, arrays & Work with categorical values and multidimensional structures. \\
5 & Selection and repetition & Select code paths and repeat operations. \\
6 & Data frames & Create, inspect, access, and modify tabular data. \\
7 & Managing data frames & Import, clean, filter, sort, and export data. \\
8 & Data visualization & Create and interpret common statistical charts. \\
9 & Correlation and scaling & Analyze relationships and transform numeric data. \\
\bottomrule
\end{tabularx}
\end{dstable}
